\documentclass[a4paper,12pt]{article}

% 页面设置
\usepackage[top=2.54cm, bottom=2.54cm, left=1.91cm, right=1.91cm]{geometry}

% 中文支持
\usepackage[UTF8]{ctex}
\usepackage{fontspec}

% 数学公式支持
\usepackage{amsmath}
\usepackage{amssymb}
\usepackage{amsfonts}
\usepackage{amsthm}

\usepackage{enumitem}

% 定理环境定义
\theoremstyle{definition}
\newtheorem*{pf}{证}
\newtheorem*{sol}{解}

% 图形支持
\usepackage{graphicx}
\usepackage{tikz}

% 超链接支持
\usepackage{hyperref}
\hypersetup{
    colorlinks=true,
    linkcolor=blue,
    filecolor=magenta,
    urlcolor=cyan,
}

% 行距设置
\linespread{1.25}

% 标题信息
\title{Mathematical Analysis II\\Week 11 Homework (May 12)}
\author{袁子强\hspace{1cm} 2025533009}
% \date{}

\begin{document}

\maketitle

\section*{Section 11.3}

1. 计算下列第二型曲线积分.

(1) $\displaystyle\int_L (x^2 + y^2)\,\mathrm{d}x + (x^2 - y^2)\,\mathrm{d}y$, $L$ 是曲线 $y = 1 - |1 - x|$ 从点 $(0,0)$ 到点 $(2,0)$;
\begin{sol}

\end{sol}
\newpage
(2) $\displaystyle\int_L \frac{\mathrm{d}x + \mathrm{d}y}{|x| + |y|}$, $L$ 是沿正方形 $A(1,0)$, $B(0,1)$, $C(-1,0)$, $D(0,-1)$ 逆时针一周的路径;
\begin{sol}

\end{sol}
\newpage
(3) $\displaystyle\int_L \frac{-x\,\mathrm{d}x + y\,\mathrm{d}y}{x^2 + y^2}$, $L$ 是圆周 $x^2 + y^2 = a^2$, 依逆时针方向一周的路径;
\begin{sol}

\end{sol}
\newpage
(4) $\displaystyle\int_L y^2\,\mathrm{d}x + xy\,\mathrm{d}y + xz\,\mathrm{d}z$, $L$ 是从 $O(0,0,0)$ 到 $A(1,0,0)$ 再到 $B(1,1,0)$ 最后到 $C(1,1,1)$ 的折线段;
\begin{sol}

\end{sol}
\newpage
(5) $\displaystyle\int_L \mathrm{e}^{x+y+z}\,\mathrm{d}x + \mathrm{e}^{x+y+z}\,\mathrm{d}y + \mathrm{e}^{x+y+z}\,\mathrm{d}z$, $L$ 是 $x = \cos\varphi$, $y = \sin\varphi$, $z = \frac{\varphi}{\pi}$ 从点 $A(1,0,0)$ 到点 $B(0,1,\frac{1}{2})$;
\begin{sol}

\end{sol}
\newpage
(6) $\displaystyle\int_L y\,\mathrm{d}x + z\,\mathrm{d}y + x\,\mathrm{d}z$, $L$ 是 $x + y = 2$ 与 $x^2 + y^2 + z^2 = 2(x + y)$ 的交线, 从原点看去是顺时针方向.
\begin{sol}

\end{sol}
\newpage
2. 求向量场 $\boldsymbol{v} = (y + z)\boldsymbol{i} + (z + x)\boldsymbol{j} + (x + y)\boldsymbol{k}$ 沿曲线 $L$ : $x = a\sin^2 t, y = 2a\sin t\cos t, z = a\cos^2 t\ (0 \leqslant t \leqslant \pi)$ 参数增加方向的曲线积分.
\begin{sol}

\end{sol}
\newpage
3. 设一质点处于弹性力场中, 弹力方向指向原点, 大小与质点离原点的距离成正比, 比例系数为 $k$, 若质点沿椭圆 $\frac{x^2}{a^2} + \frac{y^2}{b^2} = 1$ 从点 $(a, 0)$ 移到点 $(0, b)$, 求弹性力所做的功.
\begin{sol}

\end{sol}
\newpage
\end{document}